\subsection{Chuẩn Bị Dữ Liệu (Data Preparation)}


\subsubsection{Lựa chọn biến (Select Data)}

Quy trình chọn biến được thực hiện trên tập chuẩn \texttt{planet.csv} gồm 4636 dòng và 320
cột. Mục tiêu của bước này là loại các cột không mang thông tin dự báo độc lập, giảm nhiễu cho ma
trận thiết kế và tránh để mô hình học từ các biến rò rỉ target.

\paragraph{Giai đoạn 1: loại bỏ siêu dữ liệu và sai số đo lường.}
Tập \texttt{planet.csv} vẫn giữ đầy đủ hơn 300 cột theo cấu trúc NASA Exoplanet Archive.
Tuy nhiên, nhiều cột trong số này là metadata quan sát hoặc thông tin phụ của phép đo, không phải
đặc trưng vật lý trực tiếp của hành tinh. Ở bước lọc thô đầu tiên, nhóm loại các nhóm cột sau:
\begin{itemize}
    \item \textbf{Cột sai số đo lường}: các cột có hậu tố \texttt{err1}, \texttt{err2}, ví dụ
          \texttt{pl\_orbpererr1}, \texttt{pl\_radeerr2}, \texttt{st\_tefferr1}. Đây là độ bất
          định trên/dưới của phép đo, không phải giá trị vật lý trung tâm. Với bài toán hồi quy
          điểm, nhóm dùng giá trị ước lượng chính thay vì đưa cả khoảng sai số vào ma trận thiết kế.
    \item \textbf{Cột giới hạn}: các cột có hậu tố \texttt{lim}, ví dụ \texttt{pl\_orbperlim},
          \texttt{pl\_radelim}, \texttt{st\_masslim}. Các cờ này cho biết phép đo là giới hạn
          trên/dưới, dễ làm mô hình tuyến tính hiểu nhầm nếu dùng như biến số thông thường.
    \item \textbf{Cột tham chiếu và metadata phát hiện}: tên hành tinh/sao chủ, mã định danh
          \texttt{HD}, \texttt{HIP}, \texttt{TIC}, \texttt{Gaia}, phương pháp/năm/cơ sở phát hiện,
          link bài báo, kính thiên văn và thiết bị đo. Đây là thông tin mô tả quá trình con người
          quan sát dữ liệu, không trực tiếp quyết định bán kính hành tinh.
\end{itemize}

\paragraph{Giai đoạn 2: tinh chỉnh bằng domain knowledge và thống kê.}
Từ tập biến sau lọc thô, nhóm tiếp tục loại các cột có nguy cơ làm sai lệch hồi quy đa biến. Đây là
bước dùng kiến thức vật lý thiên văn để quyết định biến nào thật sự nên đi vào file
\texttt{planet.csv}.

\begin{table}[H]
\centering
\caption{Các nhóm biến bị loại}
\label{tab:part2_raw_to_planet_filter}
\resizebox{\textwidth}{!}{%
\begin{tabular}{@{}p{0.23\textwidth}p{0.24\textwidth}p{0.45\textwidth}@{}}
\hline
\textbf{Nhóm bị loại} & \textbf{Ví dụ} & \textbf{Lý do loại bỏ} \\
\hline
Rò rỉ dữ liệu khái niệm
& \texttt{pl\_radj}, \texttt{pl\_dens}, \texttt{pl\_ratror}
& \texttt{pl\_radj} là cùng bán kính hành tinh nhưng đổi đơn vị từ Trái Đất sang Sao Mộc.
\texttt{pl\_dens} chứa trực tiếp $R$ trong công thức
$\rho = M / (\frac{4}{3}\pi R^3)$, còn \texttt{pl\_ratror} là tỷ số $R_p/R_s$.
Giữ các biến này sẽ khiến mô hình suy ngược target thay vì học quan hệ dự báo độc lập. \\
\hline
Đa cộng tuyến gần hoàn hảo
& \texttt{pl\_bmassj}, \texttt{sy\_plx}
& \texttt{pl\_bmassj} chỉ là \texttt{pl\_bmasse} theo đơn vị Sao Mộc.
\texttt{sy\_plx} gần như là nghịch đảo của khoảng cách ($d \approx 1/p$).
Giữ đồng thời các cặp này làm ma trận thiết kế dễ suy biến và hệ số OLS kém ổn định. \\
\hline
Đại lượng dẫn xuất
& \texttt{st\_lum}, \texttt{st\_dens}
& Các biến này bị chi phối bởi nhiệt độ, khối lượng và bán kính sao:
$L \propto R^2T^4$ và $\rho_\star \propto M_\star/R_\star^3$. Ưu tiên biến gốc giúp mô hình dễ
diễn giải hơn và giảm trùng lặp thông tin. Riêng \texttt{st\_logg} được giữ tạm trong
\texttt{planet.csv} để pipeline đánh giá cùng các biến sao chủ, sau đó bị loại ở bước đầu pipeline. \\
\hline
Băng tần quang học
& \texttt{sy\_bmag}, \texttt{sy\_vmag}, \texttt{sy\_jmag}, \ldots
& Các độ sáng biểu kiến ở nhiều dải lọc thường tương quan mạnh với nhau và với khoảng cách quan
sát. Chúng dễ làm tăng VIF cục bộ nhưng không trực tiếp mô tả cấu trúc vật lý của hành tinh. \\
\hline
Biến hệ quy chiếu quan sát
& \texttt{pl\_tranmid}, \texttt{pl\_orbtper}, \texttt{sy\_pmra}, \texttt{sy\_pmdec}
& Bán kính hành tinh phụ thuộc vào vật chất và cấu trúc hệ sao, không phụ thuộc trực tiếp vào thời
điểm con người quan sát transit hay chuyển động riêng của hệ sao trên nền trời Trái Đất. \\
\hline
Biến thiếu quá nhiều
& \texttt{st\_age}, \texttt{st\_vsin}, \texttt{pl\_projobliq}, \texttt{pl\_trueobliq}, \texttt{pl\_occdep}
& Các thông số như tuổi sao, vận tốc quay sao hoặc độ nghiêng trục quay rất khó đo trong thực tế,
thường có tỷ lệ missing rất cao. Đưa vào MICE có thể khiến phần lớn giá trị là suy đoán và làm bóp
méo phương sai ban đầu. \\
\hline
\end{tabular}}
\end{table}

Để làm rõ hơn quyết định chọn/lọc cột, Bảng~\ref{tab:part2_column_decisions} trình bày các nhóm
cột cụ thể trong file gốc và lý do xử lý. Các cột có cùng hậu tố như \texttt{err1},
\texttt{err2}, \texttt{lim}, \texttt{reflink} được xử lý theo cùng một quy tắc vì chúng lặp lại
cho hầu hết các đại lượng vật lý.

{\scriptsize
\setlength{\tabcolsep}{3pt}
\begin{longtable}{@{}p{0.28\textwidth}p{0.12\textwidth}p{0.50\textwidth}@{}}
\caption{Quyết định xử lý chi tiết các nhóm cột từ dữ liệu chuẩn 320 cột}
\label{tab:part2_column_decisions}\\
\hline
\textbf{Cột / pattern trong dữ liệu chuẩn} & \textbf{Quyết định} & \textbf{Giải thích} \\
\hline
\endfirsthead
\caption[]{Quyết định xử lý chi tiết các nhóm cột từ dữ liệu chuẩn 320 cột (tiếp)}\\
\hline
\textbf{Cột / pattern trong dữ liệu chuẩn} & \textbf{Quyết định} & \textbf{Giải thích} \\
\hline
\endhead
\texttt{pl\_name}, \texttt{hostname}, \texttt{pl\_letter},
\texttt{hd\_name}, \texttt{hip\_name}, \texttt{tic\_id}, \texttt{gaia\_dr2\_id},
\texttt{gaia\_dr3\_id}
& Loại
& Đây là định danh hành tinh/sao chủ. Chúng không mang ý nghĩa vật lý liên tục để giải thích bán
kính. Nếu mã hóa trực tiếp, mô hình có thể học thuộc từng hệ sao thay vì học quy luật tổng quát. \\
\hline
\texttt{discoverymethod}, \texttt{disc\_year}, \texttt{disc\_refname},
\texttt{disc\_pubdate}, \texttt{disc\_locale}, \texttt{disc\_facility},
\texttt{disc\_telescope}, \texttt{disc\_instrument}
& Loại
& Đây là metadata của quá trình phát hiện. Chúng phản ánh lịch sử quan sát và năng lực thiết bị,
không phải cơ chế vật lý quyết định kích thước hành tinh. \\
\hline
\texttt{rv\_flag}, \texttt{pul\_flag}, \texttt{ptv\_flag}, \texttt{tran\_flag},
\texttt{ast\_flag}, \texttt{obm\_flag}, \texttt{micro\_flag}, \texttt{etv\_flag},
\texttt{ima\_flag}, \texttt{dkin\_flag}, \texttt{pl\_controv\_flag}, \texttt{ttv\_flag}
& Loại
& Các cờ này cho biết hành tinh được phát hiện bằng kỹ thuật nào hoặc có tranh cãi/biến thiên timing
hay không. Chúng dễ đưa thiên lệch quan sát vào mô hình, trong khi mục tiêu là dự đoán từ đại lượng
vật lý của hành tinh và sao chủ. \\
\hline
Mọi cột \texttt{*err1}, \texttt{*err2}
& Loại
& Đây là sai số trên/dưới của phép đo. Với hồi quy điểm, nhóm dùng giá trị trung tâm của đại lượng.
Đưa thêm sai số vào như feature có thể làm mô hình học chất lượng quan trắc thay vì học cấu trúc
vật lý. \\
\hline
Mọi cột \texttt{*lim}
& Loại
& Cờ giới hạn cho biết giá trị chỉ là upper/lower bound. Nếu trộn với biến vật lý liên tục, mô hình
tuyến tính dễ hiểu sai cờ đo lường này như một tín hiệu vật lý. \\
\hline
Mọi cột \texttt{*reflink}, \texttt{*\_systemref}
& Loại
& Đây là nguồn tham chiếu và hệ quy chiếu thời gian của phép đo. Chúng phục vụ truy xuất tài liệu,
không phục vụ nội suy bán kính hành tinh. \\
\hline
\texttt{sy\_snum}, \texttt{sy\_pnum}, \texttt{sy\_mnum}, \texttt{cb\_flag}
& Loại
& Số sao, số hành tinh, số mặt trăng hoặc cờ circumbinary là mô tả cấp hệ. Chúng có thể liên quan
đến bối cảnh hình thành, nhưng trong tập hiện tại tín hiệu chính cần học nằm ở quỹ đạo, transit,
khối lượng hành tinh và đặc trưng sao chủ. \\
\hline
\texttt{pl\_radj}, \texttt{pl\_dens}, \texttt{pl\_ratror}
& Loại
& Đây là nhóm rò rỉ khái niệm. \texttt{pl\_radj} là cùng target đổi đơn vị;
\texttt{pl\_dens} chứa bán kính trong công thức mật độ; \texttt{pl\_ratror} chứa trực tiếp tỷ số
bán kính hành tinh/sao. \\
\hline
\texttt{pl\_bmassj}, \texttt{sy\_plx}
& Loại
& \texttt{pl\_bmassj} trùng thông tin với \texttt{pl\_bmasse} nhưng khác đơn vị. \texttt{sy\_plx}
là thị sai, gần như nghịch đảo của \texttt{sy\_dist}. Giữ cả hai cặp sẽ làm tăng đa cộng tuyến. \\
\hline
\texttt{pl\_angsep}
& Loại
& Góc phân tách là đại lượng quan sát phụ thuộc vào khoảng cách hệ sao và hình học quan sát từ Trái
Đất. Nó không mô tả trực tiếp cấu trúc nội tại của hành tinh. \\
\hline
\texttt{pl\_orbincl}, \texttt{pl\_ratdor}, \texttt{pl\_occdep}, \texttt{pl\_rvamp}
& Loại
& Các biến này hoặc là hình học/quỹ đạo quan sát, hoặc phụ thuộc mạnh vào phương pháp đo và nhiều
trường hợp thiếu. Nhóm ưu tiên các biến transit ổn định hơn như \texttt{pl\_trandur},
\texttt{pl\_trandep}, \texttt{pl\_imppar}. \\
\hline
\texttt{pl\_tranmid}, \texttt{pl\_orbtper}, \texttt{pl\_orblper}
& Loại
& Đây là thời điểm hoặc góc pha quỹ đạo. Bán kính hành tinh không phụ thuộc trực tiếp vào thời điểm
con người ghi nhận transit hay epoch của cận điểm quỹ đạo. \\
\hline
\texttt{pl\_projobliq}, \texttt{pl\_trueobliq}, \texttt{st\_age}, \texttt{st\_vsin},
\texttt{st\_rotp}, \texttt{st\_radv}
& Loại
& Đây là các đại lượng khó đo, thường thiếu nhiều hoặc mô tả động học sao ở tầng phụ. Nếu đưa vào
MICE, phần lớn giá trị có nguy cơ là suy đoán, làm giảm độ tin cậy của phương sai và tương quan gốc. \\
\hline
\texttt{st\_spectype}, \texttt{st\_metratio}
& Loại
& \texttt{st\_spectype} là biến phân loại phức tạp cần mã hóa riêng; \texttt{st\_metratio} là thông
tin phụ cho cách biểu diễn kim loại tính. Nhóm giữ \texttt{st\_teff}, \texttt{st\_rad},
\texttt{st\_mass}, \texttt{st\_met}, \texttt{st\_logg} ở \texttt{planet.csv} để đại diện sao chủ. \\
\hline
\texttt{st\_lum}, \texttt{st\_dens}
& Loại
& Đây là đại lượng dẫn xuất từ các biến sao chủ gốc. Giữ chúng cùng \texttt{st\_teff},
\texttt{st\_rad}, \texttt{st\_mass} sẽ trùng thông tin và làm hệ số OLS khó diễn giải. \\
\hline
\texttt{rastr}, \texttt{ra}, \texttt{decstr}, \texttt{dec}, \texttt{glat}, \texttt{glon},
\texttt{elat}, \texttt{elon}, \texttt{sy\_pm}, \texttt{sy\_pmra}, \texttt{sy\_pmdec}
& Loại
& Đây là vị trí và chuyển động riêng trên bầu trời. Chúng mô tả hệ quy chiếu quan sát từ Trái Đất,
không quyết định trực tiếp bán kính hành tinh. \\
\hline
Các cột quang học \texttt{sy\_bmag}, \texttt{sy\_vmag}, \texttt{sy\_jmag},
\texttt{sy\_hmag}, \texttt{sy\_kmag}, \texttt{sy\_umag}, \texttt{sy\_gmag},
\texttt{sy\_rmag}, \texttt{sy\_imag}, \texttt{sy\_zmag}, \texttt{sy\_w1mag},
\texttt{sy\_w2mag}, \texttt{sy\_w3mag}, \texttt{sy\_w4mag}, \texttt{sy\_gaiamag},
\texttt{sy\_icmag}, \texttt{sy\_tmag}, \texttt{sy\_kepmag}
& Loại
& Các băng tần này đo độ sáng biểu kiến qua nhiều bộ lọc. Chúng tương quan mạnh với nhau và với
khoảng cách, dễ tạo nhiễu photometry và đa cộng tuyến cục bộ. \\
\hline
\texttt{pl\_nnotes}, \texttt{st\_nphot}, \texttt{st\_nrvc}, \texttt{st\_nspec},
\texttt{pl\_nespec}, \texttt{pl\_ntranspec}, \texttt{pl\_ndispec}
& Loại
& Đây là số lượng ghi chú hoặc số lượng chuỗi đo trong archive. Chúng phản ánh độ phong phú dữ liệu
quan sát, không phải cơ chế vật lý tạo nên bán kính hành tinh. \\
\hline
\end{longtable}
}

Ngược lại, \texttt{planet.csv} giữ lại 16 cột có ý nghĩa vật lý trực tiếp hơn. Bảng
\ref{tab:part2_kept_columns} tóm tắt vai trò của từng cột được giữ.

\begin{table}[H]
\centering
\caption{Các cột được giữ trong \texttt{planet.csv} và vai trò mô hình hóa}
\label{tab:part2_kept_columns}
\resizebox{\textwidth}{!}{%
\begin{tabular}{@{}p{0.20\textwidth}p{0.20\textwidth}p{0.52\textwidth}@{}}
\hline
\textbf{Cột giữ lại} & \textbf{Nhóm} & \textbf{Lý do giữ lại} \\
\hline
\texttt{pl\_orbper} & Quỹ đạo & Chu kỳ quỹ đạo phản ánh điều kiện động lực học quanh sao chủ. Được giữ ban đầu nhưng sau đó bị VIF loại khi thông tin trùng với \texttt{pl\_orbsmax}. \\
\texttt{pl\_orbsmax} & Quỹ đạo & Bán trục lớn mô tả khoảng cách quỹ đạo, ảnh hưởng đến bức xạ nhận được và điều kiện hình thành hành tinh. \\
\texttt{pl\_orbeccen} & Quỹ đạo & Độ lệch tâm mô tả hình dạng quỹ đạo, có thể liên quan đến lịch sử động lực của hệ. \\
\texttt{pl\_trandur} & Transit & Thời lượng transit chứa thông tin về hình học transit và kích thước/quỹ đạo tương đối. \\
\texttt{pl\_trandep} & Transit & Độ sâu transit liên quan đến tỷ lệ diện tích che khuất, nên có tín hiệu dự báo bán kính. \\
\texttt{pl\_imppar} & Transit & Impact parameter mô tả vị trí đường đi của hành tinh qua đĩa sao, giúp hiệu chỉnh tín hiệu transit. \\
\texttt{pl\_eqt} & Năng lượng & Nhiệt độ cân bằng đại diện môi trường nhiệt của hành tinh. Được giữ ban đầu nhưng sau đó bị VIF loại do trùng thông tin với quỹ đạo/nhiệt độ sao. \\
\texttt{pl\_insol} & Năng lượng & Thông lượng bức xạ nhận từ sao chủ, giữ lại vì phản ánh điều kiện năng lượng tác động lên hành tinh. \\
\texttt{pl\_bmasse} & Hành tinh & Khối lượng hành tinh là predictor vật lý quan trọng nhất để dự đoán bán kính. \\
\texttt{st\_teff} & Sao chủ & Nhiệt độ sao chủ ảnh hưởng đến bức xạ và môi trường quanh hành tinh. \\
\texttt{st\_rad} & Sao chủ & Bán kính sao chủ liên quan trực tiếp đến diễn giải tín hiệu transit. \\
\texttt{st\_mass} & Sao chủ & Khối lượng sao chủ được giữ tạm để pipeline đánh giá, sau đó bị loại ở bước đầu do chồng lấp thông tin. \\
\texttt{st\_met} & Sao chủ & Kim loại tính có thể liên quan đến thành phần vật chất và lịch sử hình thành hệ hành tinh. \\
\texttt{st\_logg} & Sao chủ & Trọng lực bề mặt sao được giữ tạm để pipeline đánh giá, sau đó bị loại ở bước đầu vì là đại lượng dẫn xuất/chồng lấp. \\
\texttt{sy\_dist} & Hệ sao & Khoảng cách tới hệ sao giúp mô tả thiên lệch quan sát và độ tin cậy tín hiệu. \\
\texttt{pl\_rade} & Target & Bán kính hành tinh theo đơn vị Trái Đất, là biến mục tiêu cần dự đoán. \\
\hline
\end{tabular}}
\end{table}

Sau hai giai đoạn này, \texttt{planet.csv} là bản 16 cột được rút ra từ dữ liệu chuẩn 4636 dòng,
trước khi chạy pipeline. Cụ thể, file này giữ các nhóm biến quỹ đạo, transit, năng lượng, khối
lượng hành tinh, đặc trưng sao chủ và khoảng cách hệ sao để phục vụ bài toán dự đoán
\texttt{pl\_rade}.

\subsubsection{Làm sạch dữ liệu (Clean Data)}

Pipeline tuân thủ nguyên tắc tránh data leakage: chia train/test trước, sau đó chỉ fit tham số tiền
xử lý trên train và áp dụng lại cho test.

\begin{align}
    &\text{Load data}
    \rightarrow
    \text{Domain restriction}
    \rightarrow
    \text{Train/test split}, \nonumber\\
    &\text{Fit preprocessing trên train}
    \rightarrow
    \text{Transform test bằng tham số train}.
\end{align}

\begin{lstlisting}[language=Python, caption={Tách train/test trước khi fit preprocessing}]
raw = raw[(raw["pl_rade"] <= 1.6) & (raw["pl_bmasse"] <= 10)]
train_raw, test_raw = train_test_split_frame(
    raw, test_size=0.2, random_state=RANDOM_STATE
)

pipeline = DataPipeline(target="pl_rade")
X_train, y_train = pipeline.fit_transform(train_raw)
X_test,  y_test  = pipeline.transform(test_raw, include_target=True)
\end{lstlisting}

Sau khi chia dữ liệu, train có 867 dòng và test có 217 dòng. Tất cả các bước log-transform,
Winsorization, MICE, VIF và standardization đều được fit trên train.

\subsubsection{Xây dựng dữ liệu (Construct Data)}

\paragraph{Log-transform.}
Nhiều đại lượng thiên văn có quan hệ lũy thừa, ví dụ Định luật Kepler hoặc quan hệ bức xạ. Với
quan hệ:
\begin{equation}
    y = Cx^\alpha,
\end{equation}
lấy log hai vế đưa bài toán về dạng gần tuyến tính:
\begin{equation}
    \log y = \log C + \alpha \log x.
\end{equation}

Pipeline dùng phép biến đổi:
\begin{equation}
    z = \log(1+x),
\end{equation}
và target được chuyển thành:
\begin{equation}
    y = \log(1+\texttt{pl\_rade}) = \texttt{log\_pl\_rade}.
\end{equation}

\begin{table}[H]
\centering
\caption{Ảnh hưởng của log-transform lên skewness}
\label{tab:part2_log_transform}
\resizebox{\textwidth}{!}{%
\begin{tabular}{@{}llrrrrr@{}}
\hline
\textbf{Biến gốc} & \textbf{Biến mới} & \textbf{Missing} & \textbf{Min trước} & \textbf{Max trước} &
\textbf{Skew trước} & \textbf{Skew sau} \\
\hline
\texttt{pl\_orbper} & \texttt{log\_pl\_orbper} & 0 & 0.1769 & 199.6688 & 6.8615 & 0.4901 \\
\texttt{pl\_orbsmax} & \texttt{log\_pl\_orbsmax} & 73 & 0.0050 & 0.6193 & 3.2313 & 2.5948 \\
\texttt{pl\_trandep} & \texttt{log\_pl\_trandep} & 40 & 0.0012 & 1.6536 & 10.1802 & 6.9522 \\
\texttt{pl\_insol} & \texttt{log\_pl\_insol} & 83 & 0.1440 & 8385.9000 & 3.9596 & -0.3815 \\
\texttt{pl\_bmasse} & \texttt{log\_pl\_bmasse} & 0 & 0.0364 & 10.0000 & 1.7694 & -0.2716 \\
\texttt{pl\_eqt} & \texttt{log\_pl\_eqt} & 67 & 171.7000 & 3186.0000 & 0.8788 & -0.4766 \\
\texttt{sy\_dist} & \texttt{log\_sy\_dist} & 6 & 6.8693 & 2712.5800 & 1.1829 & -1.1509 \\
\texttt{pl\_rade} & \texttt{log\_pl\_rade} & 0 & 0.3098 & 1.6000 & -0.6930 & -0.9918 \\
\hline
\end{tabular}}
\end{table}

\begin{figure}[H]
    \centering
    \includegraphics[width=\textwidth]{assets/part2/log_transform_before_after.png}
    \caption{So sánh phân phối trước và sau log-transform. Các biến có đuôi phải dài được nén
    thang đo rõ rệt, đặc biệt là \texttt{pl\_orbper}, \texttt{pl\_insol} và \texttt{pl\_bmasse}.}
    \label{fig:part2_log_transform}
\end{figure}

\paragraph{Winsorization.}
Nhóm không xóa ngoại lai mà cap giá trị ở hai đầu phân phối theo bách phân vị 1\% và 99\%:
\begin{equation}
    x' =
    \begin{cases}
        q_{0.01}, & x < q_{0.01},\\
        x,        & q_{0.01} \le x \le q_{0.99},\\
        q_{0.99}, & x > q_{0.99}.
    \end{cases}
\end{equation}

\begin{table}[H]
\centering
\caption{Tóm tắt Winsorization trên tập train}
\label{tab:part2_winsor}
\begin{tabular}{@{}lrrrrr@{}}
\hline
\textbf{Cột} & \textbf{p01} & \textbf{p99} & \textbf{Dưới p01} & \textbf{Trên p99} & \textbf{Clipped (\%)} \\
\hline
\texttt{pl\_orbeccen} & 0.0000 & 0.2994 & 0 & 8 & 1.061 \\
\texttt{pl\_trandur} & 0.6241 & 8.1895 & 9 & 9 & 2.120 \\
\texttt{pl\_imppar} & 0.0040 & 0.9395 & 9 & 9 & 2.179 \\
\texttt{st\_rad} & 0.1963 & 1.7887 & 9 & 9 & 2.076 \\
\hline
\end{tabular}
\end{table}

\begin{figure}[H]
    \centering
    \includegraphics[width=\textwidth]{assets/part2/winsorization_before_after.png}
    \caption{Boxplot trước và sau Winsorization. Các điểm cực trị được kéo về ngưỡng fit trên
    train, trong khi phần trung tâm của phân phối được giữ nguyên.}
    \label{fig:part2_winsor}
\end{figure}

\paragraph{MICE imputation.}
Sau log-transform và Winsorization, tập train vẫn còn missing values ở 9 cột. Nhóm dùng MICE
(Multiple Imputation by Chained Equations) thay cho median imputation đơn giản để bảo toàn cấu trúc
tương quan đa biến.

\begin{table}[H]
\centering
\caption{Missing values sau log-transform trên tập train}
\label{tab:part2_missing_after_transform}
\begin{tabular}{@{}lrrl@{}}
\hline
\textbf{Cột} & \textbf{Missing count} & \textbf{Missing (\%)} & \textbf{Kiểu dữ liệu} \\
\hline
\texttt{pl\_orbeccen} & 113 & 13.033 & float64 \\
\texttt{log\_pl\_insol} & 83 & 9.573 & float64 \\
\texttt{log\_pl\_orbsmax} & 73 & 8.420 & float64 \\
\texttt{log\_pl\_eqt} & 67 & 7.728 & float64 \\
\texttt{pl\_imppar} & 41 & 4.729 & float64 \\
\texttt{log\_pl\_trandep} & 40 & 4.614 & float64 \\
\texttt{st\_met} & 24 & 2.768 & float64 \\
\texttt{pl\_trandur} & 18 & 2.076 & float64 \\
\texttt{log\_sy\_dist} & 6 & 0.692 & float64 \\
\hline
\end{tabular}
\end{table}

MICE trong pipeline dùng 5 bản imputation, 5 vòng chained iterations, hồi quy Ridge phụ với
regularization 0.001 và noise scale 0.15.

\begin{table}[H]
\centering
\caption{Tóm tắt giá trị được MICE impute trên tập train}
\label{tab:part2_mice}
\resizebox{\textwidth}{!}{%
\begin{tabular}{@{}lrrrrrr@{}}
\hline
\textbf{Cột} & \textbf{Missing} & \textbf{Missing (\%)} & \textbf{Median khởi tạo} &
\textbf{Observed mean} & \textbf{Imputed mean} & \textbf{Imputed std} \\
\hline
\texttt{pl\_orbeccen} & 113 & 13.033 & 0.000000 & 0.013696 & 0.033514 & 0.021444 \\
\texttt{log\_pl\_insol} & 83 & 9.573 & 5.207544 & 5.036231 & 5.332386 & 1.555943 \\
\texttt{log\_pl\_orbsmax} & 73 & 8.420 & 0.052213 & 0.064713 & 0.045031 & 0.040655 \\
\texttt{log\_pl\_eqt} & 67 & 7.728 & 6.856987 & 6.806412 & 6.805654 & 0.457823 \\
\texttt{pl\_imppar} & 41 & 4.729 & 0.284700 & 0.346475 & 0.394104 & 0.059489 \\
\texttt{log\_pl\_trandep} & 40 & 4.614 & 0.020498 & 0.036569 & 0.046997 & 0.040054 \\
\texttt{st\_met} & 24 & 2.768 & 0.000000 & -0.025205 & -0.097143 & 0.039923 \\
\texttt{pl\_trandur} & 18 & 2.076 & 2.372000 & 2.699806 & 2.127761 & 1.064210 \\
\texttt{log\_sy\_dist} & 6 & 0.692 & 6.104394 & 5.830641 & 5.143891 & 0.991066 \\
\hline
\end{tabular}}
\end{table}

\begin{figure}[H]
    \centering
    \includegraphics[width=\textwidth]{assets/part2/mice_observed_vs_imputed.png}
    \caption{So sánh phân phối observed và phần được MICE impute. Giá trị impute không bị gom
    về một điểm trung vị duy nhất mà vẫn có độ phân tán theo quan hệ đa biến.}
    \label{fig:part2_mice}
\end{figure}

\paragraph{Giai đoạn 3 trong DataPipeline: loại biến bằng VIF.}
Khác với hai giai đoạn tạo \texttt{planet.csv} ở trên, bước VIF diễn ra bên trong
\texttt{DataPipeline}. Sau khi train/test đã được tách và tập train đã đi qua log-transform,
Winsorization và MICE, pipeline kiểm tra đa cộng tuyến bằng hệ số phóng đại phương sai:
\begin{equation}
    \mathrm{VIF}_j = \frac{1}{1-R_j^2},
\end{equation}
trong đó $R_j^2$ là hệ số xác định khi hồi quy biến $x_j$ theo các biến còn lại. Ngưỡng loại biến
là $\mathrm{VIF} > 10$.

\begin{table}[H]
\centering
\caption{Lịch sử lọc VIF}
\label{tab:part2_vif_history}
\begin{tabular}{@{}lrl@{}}
\hline
\textbf{Biến có VIF cao nhất} & \textbf{VIF} & \textbf{Quyết định} \\
\hline
\texttt{log\_pl\_orbper} & 43.1374 & Drop \\
\texttt{log\_pl\_eqt} & 20.4514 & Drop \\
\texttt{st\_teff} & 9.2570 & Dừng vì nhỏ hơn 10 \\
\hline
\end{tabular}
\end{table}

\begin{table}[H]
\centering
\caption{VIF cuối cùng của các feature còn lại}
\label{tab:part2_final_vif}
\begin{tabular}{@{}lrlr@{}}
\hline
\textbf{Feature} & \textbf{VIF} & \textbf{Feature} & \textbf{VIF} \\
\hline
\texttt{st\_teff} & 9.2570 & \texttt{log\_pl\_trandep} & 1.8730 \\
\texttt{st\_rad} & 5.9915 & \texttt{pl\_orbeccen} & 1.2736 \\
\texttt{log\_pl\_insol} & 5.8778 & \texttt{st\_met} & 1.1861 \\
\texttt{log\_pl\_orbsmax} & 4.6054 & \texttt{pl\_imppar} & 1.0732 \\
\texttt{pl\_trandur} & 3.0931 & \texttt{log\_pl\_bmasse} & 1.0598 \\
\texttt{log\_sy\_dist} & 2.8147 & & \\
\hline
\end{tabular}
\end{table}

\begin{figure}[H]
    \centering
    \includegraphics[width=\textwidth]{assets/part2/vif_before_after.png}
    \caption{VIF trước và sau lọc. Việc loại \texttt{log\_pl\_orbper} và \texttt{log\_pl\_eqt}
    đưa tất cả VIF còn lại xuống dưới ngưỡng 10.}
    \label{fig:part2_vif}
\end{figure}

Sau VIF filtering, 11 feature cuối cùng là:
\begin{center}
\begin{tabular}{@{}lll@{}}
\hline
\texttt{pl\_orbeccen} & \texttt{pl\_trandur} & \texttt{pl\_imppar} \\
\texttt{st\_teff} & \texttt{st\_rad} & \texttt{st\_met} \\
\texttt{log\_pl\_orbsmax} & \texttt{log\_pl\_trandep} & \texttt{log\_pl\_insol} \\
\texttt{log\_pl\_bmasse} & \texttt{log\_sy\_dist} & \\
\hline
\end{tabular}
\end{center}

\paragraph{Chuẩn hóa dữ liệu.}
Tất cả feature cuối cùng được chuẩn hóa bằng z-score:
\begin{equation}
    x_{ij}^{\mathrm{std}} =
    \frac{x_{ij} - \mu_j^{\mathrm{train}}}{\sigma_j^{\mathrm{train}}},
\end{equation}
trong đó $\mu_j^{\mathrm{train}}$ và $\sigma_j^{\mathrm{train}}$ chỉ được fit trên train.


